\ifdefined\gkver\else\def\gkver{student}\fi
\documentclass[\gkver]{gaokaozhenti}
\begin{document}

\chapter{2016年海南}

\begin{enumerate}
\item
%% number: 1
%% paperName: 2016年海南高考第 1 题 3 分
%% typeId: 1
%% type: 单选题
%% chapter: 曲线运动
%% point: 平抛运动
%% method: 无
%% score: 3
%% degree: 750
%% duplicateId: 0
%% body:
在地面上方某点将一小球以一定的初速度沿水平方向抛出，不计空气阻力，则小球在随后的运动中，\xzanswer{B}
\fourchoices[answer=B]
{速度和加速度的方向都在不断变化}
{速度与加速度方向之间的夹角一直减小}
{在相等的时间间隔内，速率的改变量相等}
{在相等的时间间隔内，动能的改变量相等}
%% answer: B
%% memo:
\memoanswer{
做平抛运动的物体加速度不变，速度不断改变，故 A 错误；做平抛运动的物体的运动方向向着重力的方向偏转，故速度与加速度方向间的夹角一直减小，故 B 正确；在相等时间间隔内，速度改变量相等，但速度是矢量，速率为标量，故速率改变量不等，故 C 错误；相等时间间隔内竖直位移不断增大，重力做功不相等，由动能定理，动能的改变量不相等，故 D 错误。

\twopicture[align=center, wA=4.2cm, wB=4.6cm]{figs/01_an1.png}{figs/01_an2.png}
}

\item
%% number: 2
%% paperName: 2016年海南高考第 2 题 3 分
%% typeId: 1
%% type: 单选题
%% chapter: 力物体的平衡
%% point: 物体系平衡
%% method: 整体与隔离
%% score: 3
%% degree: 550
%% duplicateId: 0
%% body:
如图，在水平桌面上放置一斜面体 P，两长方体物块 a 和 b 叠放在 P 的斜面上，整个系统处于静止状态。若将 a 和 b、b 与 P、P 与桌面之间摩擦力的大小分别用 $f_{1}$、$f_{2}$ 和 $f_{3}$ 表示。则\xzanswer{C}
\onepicture[w=5.2cm]{figs/02.png}
\fourchoices[answer=C]
{$f_{1}=0$，$f_{2}\neq0$，$f_{3}\neq0$}
{$f_{1}\neq0$，$f_{2}=0$，$f_{3}=0$}
{$f_{1}\neq0$，$f_{2}\neq0$，$f_{3}=0$}
{$f_{1}\neq0$，$f_{2}\neq0$，$f_{3}\neq0$}
%% answer: C
%% memo:
\memoanswer{
利用整体法分析，水平方向上不受力，故 P 与桌面间的摩擦力 $f_{3}=0$，研究整体 ab 受重力、支持力和摩擦力，故 $f_{2}\neq0$，研究 a 受重力、支持力和摩擦力，故为 $f_{1}\neq0$，故 C 正确。
}

\item
%% number: 3
%% paperName: 2016年海南高考第 3 题 3 分
%% typeId: 1
%% type: 单选题
%% chapter: 曲线运动
%% point: 竖直平面圆周运动
%% method: 无
%% score: 3
%% degree: 450
%% duplicateId: 0
%% body:
如图，光滑圆轨道固定在竖直面内，一质量为 $m$ 的小球沿轨道做完整的圆周运动。已知小球在最低点时对轨道的压力大小为 $N_{1}$，在高点时对轨道的压力大小为 $N_{2}$。重力加速度大小为 $g$，则 $N_{1}-N_{2}$ 的值为\xzanswer{D}
\onepicture[w=3cm]{\includesvg{figs/03}}
\fourchoices[answer=D]
{$3mg$}
{$4mg$}
{$5mg$}
{$6mg$}
%% answer: D
%% memo:
\memoanswer{
在最低点 $N_{1}-mg=m\frac{v_{1}^{2}}{R}$，在最高点 $N_{2}+mg=m\frac{v_{2}^{2}}{R}$，小球运动过程中机械能守恒，故 $\frac{1}{2}mv_{1}^{2}=mg\cdot2R+\frac{1}{2}mv_{2}^{2}$，解得 $N_{1}-N_{2}=6mg$，故 D 正确。
}

\item
%% number: 4
%% paperName: 2016年海南高考第 4 题 3 分
%% typeId: 1
%% type: 单选题
%% chapter: 电磁感应
%% point: 楞次定律
%% method: 无
%% score: 3
%% degree: 550
%% duplicateId: 0
%% body:
如图，一圆形金属环与两固定的平行长直导线在同一竖直平面内，环的圆心与两导线距离相等，环的直径小于两导线间距。两导线中通有大小相等、方向向下的恒定电流。若\xzanswer{D}
\onepicture[w=2.4cm]{\includesvg{figs/04}}
\fourchoices[answer=D]
{金属环向上运动，则环上的感应电流方向为顺时针方向}
{金属环向下运动，则环上的感应电流方向为顺时针方向}
{金属环向左侧直导线靠近，则环上的感应电流方向为逆时针方向}
{金属环向右侧直导线靠近，则环上的感应电流方向为逆时针方向}
%% answer: D
%% memo:
\memoanswer{
若金属环向上或向下运动，则由对称性，金属环中的磁通量为 0 且不发生改变，金属环中没有感应电流，故 AB 错误；金属环向左侧靠近导线，则环内向外的磁通量增加，由楞次定律，感应电流产生的磁场方向向里，由安培定则知，感应电流的方向为顺时针方向，故 C 错误；金属环向右侧靠近导线，则环内向里的磁通量增加，由楞次定律，感应电流的磁场方向向外，由安培定则知，感应电流的方向为逆时针方向，故 D 正确。
}

\item
%% number: 5
%% paperName: 2016年海南高考第 5 题 3 分
%% typeId: 1
%% type: 单选题
%% chapter: 运动和力
%% point: 牛顿定律中的图象
%% method: 无
%% score: 3
%% degree: 550
%% duplicateId: 0
%% body:
沿固定斜面下滑的物体受到与斜面平行向上的拉力 $F$ 的作用，其下滑的速度-时间图线如图所示。已知物体与斜面之间的动摩擦因数为常数，在 $0\sim5\Us$、$5\sim10\Us$、$10\sim15\Us$ 内 $F$ 的大小分别为 $F_{1}$、$F_{2}$ 和 $F_{3}$，则\xzanswer{A}
\onepicture[w=5cm]{\includesvg{figs/05}}
\fourchoices[answer=A]
{$F_{1}<F_{2}$}
{$F_{2}>F_{3}$}
{$F_{1}>F_{3}$}
{$F_{1}=F_{3}$}
%% answer: A
%% memo:
\memoanswer{
加速下滑的过程 $mg\sin\theta-F_{1}-f=ma$，匀速下滑的过程 $mg\sin\theta-F_{2}-f=0$，减速下滑的过程 $F_{3}-mg\sin\theta+f=ma$，故 $F_{1}<F_{2}<F_{3}$，故 A 正确。
}

\item
%% number: 6
%% paperName: 2016年海南高考第 6 题 3 分
%% typeId: 1
%% type: 单选题
%% chapter: 电场
%% point: 带电粒子的偏转
%% method: 无
%% score: 3
%% degree: 350
%% duplicateId: 0
%% body:
如图，平行板电容器两极板的间距为 $d$，极板与水平面成 $45^{\circ}$ 角，上极板带正电。一电荷量为 $q$（$q>0$）的粒子在电容器中靠近下极板处，以初动能 $E_{k0}$ 竖直向上射出。不计重力，极板尺寸足够大。若粒子能打到上极板，则两极板间电场强度的最大值为\xzanswer{B}
\onepicture[w=3cm]{\includesvg{figs/06}}
\fourchoices[answer=B]
{$\frac{E_{k0}}{4qd}$}
{$\frac{E_{k0}}{2qd}$}
{$\frac{\sqrt{2}E_{k0}}{2qd}$}
{$\frac{\sqrt{2}E_{k0}}{qd}$}
%% answer: B
%% memo:
\memoanswer{
若粒子刚好能打到上极板，则垂直于平行板的速度刚好减为 0，因极板与水平面夹角为 $45^{\circ}$，末速度为初速度的 $\frac{\sqrt{2}}{2}$，故粒子的末动能为 $\frac{1}{2}E_{k0}$，由动能定理得 $-qEd=\frac{1}{2}E_{k0}-E_{k0}$ 可得 $E=\frac{E_{k0}}{2qd}$，故 B 正确。
}

\item
%% number: 7
%% paperName: 2016年海南高考第 7 题 5 分
%% typeId: 2
%% type: 多选题
%% chapter: 曲线运动
%% point: 人造地球卫星
%% method: 无
%% score: 5
%% degree: 450
%% duplicateId: 0
%% body:
通过观测冥王星的卫星，可以推算出冥王星的质量。假设卫星绕冥王星做匀速圆周运动，除了引力常量外，至少还需要两个物理量才能计算出冥王星的质量。这两个物理量可以是\xzanswer{AD}
\fourchoices[answer=AD]
{卫星的速度和角速度}
{卫星的质量和轨道半径}
{卫星的质量和角速度}
{卫星的运行周期和轨道半径}
%% answer: AD
%% memo:
\memoanswer{
知道卫星的速度和角速度，由 $v=\omega r$ 可求得卫星的轨道半径，由 $G\frac{Mm}{r^{2}}=m\frac{v^{2}}{r}$，即可求得冥王星的质量，故 A 正确；由 $G\frac{Mm}{r^{2}}=m\left(\frac{2\pi}{T}\right)^{2}r$，知道卫星的运行周期和轨道半径，可求得冥王星的质量，故 D 正确；求冥王星的质量，不需要知道卫星的质量，故 BC 错误。
}

\item
%% number: 8
%% paperName: 2016年海南高考第 8 题 5 分
%% typeId: 2
%% type: 多选题
%% chapter: 磁场
%% point: 左手定则
%% method: 微元
%% score: 5
%% degree: 550
%% duplicateId: 0
%% body:
如图\ref{2016海南08a}所示，扬声器中有一线圈处于磁场中，当音频电流信号通过线圈时，线圈带动纸盆振动，发出声音。俯视图\ref{2016海南08b}表示处于辐射状磁场中的线圈（线圈平面即纸面），磁场方向如图中箭头所示。在图\ref{2016海南08b}中\xzanswer{BC}
\twopicture[labelA=2016海南08a, labelB=2016海南08b, align=center, wA=7cm, wB=3.6cm]
{figs/08a.png}
{figs/08b.png}
\fourchoices[answer=BC]
{当电流沿顺时针方向时，线圈所受安培力的方向垂直于纸面向里}
{当电流沿顺时针方向时，线圈所受安培力的方向垂直于纸面向外}
{当电流沿逆时针方向时，线圈所受安培力的方向垂直于纸面向里}
{当电流沿逆时针方向时，线圈所受安培力的方向垂直于纸面向外}
%% answer: BC
%% memo:
\memoanswer{
由安培定则，当电流沿顺时针方向时，线圈所受安培力的方向垂直纸面向外，故 A 错误 B 正确；当电流沿逆时针方向时，线圈所受安培力的方向垂直纸面向里，故 C 正确 D 错误。
}

\item
%% number: 9
%% paperName: 2016年海南高考第 9 题 5 分
%% typeId: 2
%% type: 多选题
%% chapter: 交流电
%% point: 变压器
%% method: 无
%% score: 5
%% degree: 550
%% duplicateId: 0
%% body:
图\ref{2016海南09a}所示，理想变压器的原、副线圈的匝数比为 $4:1$，$R_{T}$ 为阻值随温度升高而减小的热敏电阻，$R_{1}$ 为定值电阻，电压表和电流表均为理想交流电表。原线圈所接电压 $u$ 随时间 $t$ 按正弦规律变化，如图\ref{2016海南09b}所示。下列说法正确的是\xzanswer{BD}
\twopicture[labelA=2016海南09a, labelB=2016海南09b, align=center, wA=4.6cm, wB=5cm]
{figs/09a.png}
{figs/09b.png}
\fourchoices[answer=BD]
{变压器输入、输出功率之比为 $4:1$}
{变压器原、副线圈中的电流强度之比为 $1:4$}
{$u$ 随 $t$ 变化的规律为 $u=51\sin(50\pi t)$（国际单位制）}
{若热敏电阻 $R_{T}$ 的温度升高，则电压表的示数不变，电流表的示数变大}
%% answer: BD
%% memo:
\memoanswer{
理想变压器输入、输出功率相等，故 A 错误；由 $\frac{I_{1}}{I_{2}}=\frac{n_{2}}{n_{1}}=\frac{1}{4}$，故 B 正确；由图象，$u$ 随 $t$ 变化的规律为 $u=51\sin(100\pi t)\UV$，故 C 错误；若热敏电阻温度升高，则副线圈电路中电阻减小，原、副线圈电压不变，电流表示数变大，故 D 正确。
}

\item
%% number: 10
%% paperName: 2016年海南高考第 10 题 5 分
%% typeId: 2
%% type: 多选题
%% chapter: 电场
%% point: 带电粒子的轨迹类问题
%% method: 无
%% score: 5
%% degree: 450
%% duplicateId: 0
%% body:
如图，一带正电的点电荷固定于 O 点，两虚线圆均以 O 为圆心，两实线分别为带电粒子 M 和 N 先后在电场中运动的轨迹，a、b、c、d、e 为轨迹和虚线圆的交点。不计重力。下列说法说法正确的是\xzanswer{ABC}
\onepicture[w=4.6cm]{\includesvg{figs/10}}
\fourchoices[answer=ABC]
{M 带负电荷，N 带正电荷}
{M 在 b 点的动能小于它在 a 点的动能}
{N 在 d 点的电势能等于它在 e 点的电势能}
{N 在从 c 点运动到 d 点的过程中克服电场力做功}
%% answer: ABC
%% memo:
\memoanswer{
做曲线运动的物体一定受到指向轨迹内侧的合外力，故 M 带负电荷、N 带正电荷，故 A 正确；对 M 由 $W_{ab}=(-q)U_{ab}<0$，M 从 a 到 b 电场力做负功，动能减小，故 B 正确；d、e 在同一等势面，N 在 d 点的电势能等于它在 e 点的电势能，故 C 正确；对 N 由 $W_{cd}=(+q)U_{cd}>0$，即 N 在从 c 点运动到 d 点的过程中电场力做正功，故 D 错误。
}

\item
%% number: 11
%% paperName: 2016年海南高考第 11 题 6 分
%% typeId: 4
%% type: 实验题
%% chapter: 直线运动
%% point: 打点计时器公式
%% method: 无
%% score: 6
%% degree: 550
%% duplicateId: 0
%% body:
某同学利用图\ref{2016海南11a}所示的实验装置探究物块速度随时间的变化。物块放在桌面上，细绳的一端与物块相连，另一端跨过滑轮挂上钩码。打点计时器固定在桌面左端，所用交流电源频率为 $50\UHz$。纸带穿过打点计时器连接在物块上。启动打点计时器，释放物块，物块在钩码的作用下拖着纸带运动。打点计时器打出的纸带如图\ref{2016海南11b}所示（图中相邻两点间有 4 个点未画出）。

根据实验数据分析，该同学认为物块的运动为匀加速运动。回答下列问题：
\begin{enumerate}
	\item
	在打点计时器打出 B 点时，物块的速度大小为\tkanswer[2cm]{0.56}$\Ums$。在打出 D 点时，物块的速度大小为\tkanswer[2cm]{0.96}$\Ums$；（保留两位有效数字）
	\item
	物块的加速度大小为\tkanswer[2cm]{2.0}$\Umsq$。（保留两位有效数字）
\end{enumerate}
\twopicture[labelA=2016海南11a, labelB=2016海南11b, align=right, wA=7cm, wB=9cm]
{figs/11a.png}
{figs/11b.png}
%% answer:
\jdanswer{
\begin{enumerate}
	\item
	$0.56\Ums$，$0.96\Ums$。
	\item
	$2.0\Umsq$。
\end{enumerate}
}
%% memo:
\memoanswer{
（1）由匀变速运动中间时刻的瞬时速度等于整段的平均速度，可得 $v_{B}=\frac{(4.61+6.59)\times10^{-2}}{2\times5\times0.02}\Ums=0.56\Ums$，$v_{D}=\frac{(8.61+10.61)\times10^{-2}}{2\times5\times0.02}\Ums=0.96\Ums$。

（2）由 $a=\frac{\Delta v}{\Delta t}$ 可得 $a=\frac{0.96-0.56}{2\times5\times0.02}\Umsq=2.0\Umsq$。
}

\item
%% number: 12
%% paperName: 2016年海南高考第 12 题 9 分
%% typeId: 4
%% type: 实验题
%% chapter: 电路
%% point: 电路实验
%% method: 无
%% score: 9
%% degree: 450
%% duplicateId: 0
%% body:
某同学改装和校准电压表的电路图如图所示，图中虚线框内是电压表的改装电路。
\begin{enumerate}
	\item
	已知表头 G 满偏电流为 $100\UmuA$，表头上标记的内阻值为 $900\UO$。$R_{1}$、$R_{2}$ 和 $R_{3}$ 是定值电阻。利用 $R_{1}$ 和表头构成 $1\UmA$ 的电流表，然后再将其改装为两个量程的电压表。若使用 a、b 两个接线柱，电压表的量程为 $1\UV$；若使用 a、c 两个接线柱，电压表的量程为 $3\UV$。则根据题给条件，定值电阻的阻值应选 $R_{1}=$\tkanswer[2cm]{100}$\UO$，$R_{2}=$\tkanswer[2cm]{910}$\UO$，$R_{3}=$\tkanswer[2cm]{2000}$\UO$。
	\item
	用量程为 $3\UV$，内阻为 $2500\UO$ 的标准电压表 V 对改装表 $3\UV$ 挡的不同刻度进行校准。所用电池的电动势 $E$ 为 $5\UV$；滑动变阻器 $R$ 有两种规格，最大阻值分别为 $50\UO$ 和 $5\UkO$。为方便实验中调节电压，图中 $R$ 应选用最大阻值为\tkanswer[2cm]{50}$\UO$ 的滑动变阻器。
	\item
	校准时，在闭合开关 S 前，滑动变阻器的滑动端 P 应靠近\tkanswer[1.5cm]{M}（填“M”或“N”）端。
	\item
	若由于表头 G 上标记的内阻值不准，造成改装后电压表的读数比标准电压表的读数偏小，则表头 G 内阻的真实值\tkanswer[1.5cm]{大于}（填“大于”或“小于”）$900\UO$。
\end{enumerate}
\onepicture[w=5.5cm, align=right]{figs/12.png}
%% answer:
\jdanswer{
\begin{enumerate}
	\item
	$100\UO$，$910\UO$，$2000\UO$。
	\item
	$50\UO$。
	\item
	M。
	\item
	大于。
\end{enumerate}
}
%% memo:
\memoanswer{
（1）由并联分流 $I_{g}r_{g}=\left(I-I_{g}\right)R_{1}$，解得 $R_{1}=100\UO$；$R_{1}$ 与表头的并联电阻 $R_{\text{并}}=90\UO$，$U_{g}=I_{g}r_{g}=0.09\UV$，由串联分压 $\frac{U_{g}}{R_{\text{并}}}=\frac{U-U_{g}}{R}$，当 $U=1\UV$ 时，解得 $R_{2}=910\UO$，当 $U=3\UV$ 时，$R_{2}+R_{3}=2910\UO$，解得 $R_{3}=2000\UO$。

（2）滑动变阻器用分压式接法时，应选择阻值小的，此时测量电路的电压接近线性变化，易于调节，故选择 $50\UO$ 的滑动变阻器。

（3）为了保证闭合开关时，测量电路部分电压最小，故闭合开关前滑片滑到 M 端。

（4）电压表读数偏小，即相等电压情况下，电流计偏转比标准电压表小，并联电路中电阻越大，电流越小，故电流计的实际电阻偏大。
}

\item
%% number: 13
%% paperName: 2016年海南高考第 13 题 9 分
%% typeId: 6
%% type: 计算题
%% chapter: 运动和力
%% point: 牛顿定律的应用
%% method: 无
%% score: 9
%% degree: 650
%% duplicateId: 0
%% body:
水平地面上有质量分别为 $m$ 和 $4m$ 的物块 A 和 B，两者与地面的动摩擦因数均为 $\mu$。细绳的一端固定，另一端跨过轻质动滑轮与 A 相连，动滑轮与 B 相连，如图所示。初始时，绳处于水平拉直状态。若物块 A 在水平向右的恒力 $F$ 作用下向右移动了距离 $s$，重力加速度大小为 $g$。求：
\begin{enumerate}
	\item
	物块 B 克服摩擦力所做的功；
	\item
	物块 A、B 的加速度大小。
\end{enumerate}
\onepicture[w=7cm, align=right]{figs/13.png}
%% answer:
\jdanswer{
\begin{enumerate}
	\item
	$2\mu mgs$。
	\item
	$\frac{F-3\mu mg}{2m}$，$\frac{F-3\mu mg}{4m}$。
\end{enumerate}
}
%% memo:
\memoanswer{
（1）物块 A 移动了距离 $s$，则物块 B 移动的距离为 $s_{1}=\frac{1}{2}s$①，物块 B 受到的摩擦力大小为 $f=4\mu mg$②，物块 B 克服摩擦力所做的功为 $W=fs_{1}=2\mu mgs$③。

（2）设物块 A、B 的加速度大小分别为 $a_{A}$、$a_{B}$，绳中的张力为 $T$。由牛顿第二定律得 $F-\mu mg-T=ma_{A}$④，$2T-4\mu mg=ma_{B}$⑤，由 A 和 B 的位移关系得 $a_{A}=2a_{B}$⑥，联立④⑤⑥式得 $a_{A}=\frac{F-3\mu mg}{2m}$⑦，$a_{B}=\frac{F-3\mu mg}{4m}$⑧。
}

\item
%% number: 14
%% paperName: 2016年海南高考第 14 题 14 分
%% typeId: 6
%% type: 计算题
%% chapter: 磁场
%% point: 带电粒子在磁场中的运动
%% method: 无
%% score: 14
%% degree: 350
%% duplicateId: 0
%% body:
如图，A、C 两点分别位于 $x$ 轴和 $y$ 轴上，$\angle OCA=30^{\circ}$，OA 的长度为 $L$。在 $\triangle OCA$ 区域内有垂直于 $xOy$ 平面向里的匀强磁场。质量为 $m$、电荷量为 $q$ 的带正电粒子，以平行于 $y$ 轴的方向从 OA 边射入磁场。已知粒子从某点射入时，恰好垂直于 OC 边射出磁场，且粒子在磁场中运动的时间为 $t_{0}$。不计重力。
\begin{enumerate}
	\item
	求磁场的磁感应强度的大小；
	\item
	若粒子先后从两不同点以相同的速度射入磁场，恰好从 OC 边上的同一点射出磁场，求该粒子这两次在磁场中运动的时间之和；
	\item
	若粒子从某点射入磁场后，其运动轨迹与 AC 边相切，且在磁场内运动的时间为 $\frac{5}{3}t_{0}$，求粒子此次入射速度的大小。
\end{enumerate}
\onepicture[w=4.2cm, align=right]{figs/14.png}
%% answer:
\jdanswer{
\begin{enumerate}
	\item
	$\frac{\pi m}{2qt_{0}}$。
	\item
	$2t_{0}$。
	\item
	$\frac{\sqrt{3}\pi L}{7t_{0}}$。
\end{enumerate}
}
%% memo:
\memoanswer{
（1）粒子在磁场中做匀速圆周运动，在时间 $t_{0}$ 内其速度方向改变了 $90^{\circ}$，故其周期 $T=4t_{0}$①，设磁感应强度大小为 $B$，粒子的速度大小为 $v$，圆周运动的半径为 $r$，由洛伦兹力公式和牛顿第二定律得 $qvB=m\frac{v^{2}}{r}$②，匀速圆周运动速度满足 $v=\frac{2\pi r}{T}$③，联立①②③式得 $B=\frac{\pi m}{2qt_{0}}$④。

（2）设粒子从 OA 边两个不同位置射入磁场，能从 OC 边上的同一点 P 射出磁场，粒子在磁场中运动的轨迹如图（a）所示。设两轨迹所对应的圆心角分别为 $\theta_{1}$ 和 $\theta_{2}$，由几何关系得 $\theta_{1}=180^{\circ}-\theta_{2}$⑤，粒子两次在磁场中运动的时间分别为 $t_{1}$ 与 $t_{2}$，则 $t_{1}+t_{2}=\frac{T}{2}=2t_{0}$⑥。

\onepicture[w=3.8cm, align=right]{figs/14_an1.png}

（3）如图（b），由题给条件可知，该粒子在磁场区域中的轨迹圆弧对应的圆心角为 $150^{\circ}$。设 $O'$ 为圆弧的圆心，圆弧的半径为 $r_{0}$，圆弧与 AC 相切于 B 点，从 D 点射出磁场，由几何关系和题给条件可知，此时有 $\angle OO'D=\angle BO'A=30^{\circ}$⑦，$r_{0}\cos\angle OO'D+\frac{r_{0}}{\cos\angle BO'A}=L$⑧，设粒子此次入射速度的大小为 $v_{0}$，由圆周运动规律得 $v_{0}=\frac{2\pi r_{0}}{T}$⑨。

\onepicture[w=4.2cm, align=right]{figs/14_an2.png}

联立①⑦⑧⑨式得 $v_{0}=\frac{\sqrt{3}\pi L}{7t_{0}}$⑩。
}

\item
%% number: 15
%% paperName: 2016年海南高考第 15 题 8 分
%% typeId: 6
%% type: 计算题
%% chapter: 气体性质
%% point: 气体连接体
%% method: 无
%% score: 8
%% degree: 450
%% duplicateId: 0
%% body:
{}[选修 3-3]（12 分）
\begin{enumerate}
	\item
	（4 分）一定量的理想气体从状态 M 可以经历过程 1 或者过程 2 到达状态 N，其 $p-V$ 图象如图所示。在过程 1 中，气体始终与外界无热量交换；在过程 2 中，气体先经历等容变化再经历等压变化。对于这两个过程，下列说法正确的是\tkanswer[1.6cm]{ABE}。（填入正确答案标号。选对 1 个得 2 分，选对 2 个得 3 分，选对 3 个得 4 分；有选错的得 0 分）
	\onepicture[w=5.2cm, align=right]{figs/15a.png}
	\fivechoices[answer=ABE]
	{气体经历过程 1，其温度降低}
	{气体经历过程 1，其内能减小}
	{气体在过程 2 中一直对外放热}
	{气体在过程 2 中一直对外做功}
	{气体经历过程 1 的内能改变量与经历过程 2 的相同}
	\item
	（8 分）如图，密闭汽缸两侧与一 U 形管的两端相连，汽缸壁导热；U 形管内盛有密度为 $\rho=7.5\times10^{2}\Ukgmc$ 的液体。一活塞将汽缸分成左、右两个气室，开始时，左气室的体积是右气室的体积的一半，气体的压强均为 $p_{0}=4.5\times10^{3}\UPa$。外界温度保持不变。缓慢向右拉活塞使 U 形管两侧液面的高度差 $h=40\Ucm$，求此时左、右两气室的体积之比。取重力加速度大小 $g=10\Umsq$，U 形管中气体的体积和活塞拉杆的体积忽略不计。
	\onepicture[w=4.5cm, align=right]{figs/15b.png}
\end{enumerate}
%% answer:
\jdanswer{
\begin{enumerate}
	\item
	ABE。
	\item
	$1:1$。
\end{enumerate}
}
%% memo:
\memoanswer{
（1）过程 1 气体与外界没有热量交换，体积增大，气体对外界做功，由热力学第一定律，内能必定减小，温度降低，故 AB 正确；气体在过程 2 中，第一过程为等容变化，压强减小，温度降低，内能减小，对外放热，没有对外做功，第二过程为等压变化，体积增大，温度升高，内能增大，但对外做功，由热力学第一定律，一定从外界吸热，故 CD 错误；内能的改变量只与初末状态有关，与过程无关，两个过程内能改变量相同，故 E 正确。

（2）设初始状态时气缸左气室的体积为 $V_{01}$，右气室的体积为 $V_{02}$；当活塞至气缸中某位置时，左、右气室的压强分别为 $p_{1}$、$p_{2}$，体积分别为 $V_{1}$、$V_{2}$，由玻意耳定律得 $p_{0}V_{01}=p_{1}V_{1}$①，$p_{0}V_{02}=p_{2}V_{2}$②，依题意有 $V_{01}+V_{02}=V_{1}+V_{2}$③，由力的平衡条件得 $p_{2}-p_{1}=\rho gh$④，联立①②③④式，并代入题给数据得 $2V_{1}^{2}+3V_{01}V_{1}-9V_{01}^{2}=0$⑤，由此解得 $V_{1}=\frac{3}{2}V_{01}$（另一解不合题意，舍去）⑥，由③⑥式和题给条件得 $V_{1}:V_{2}=1:1$⑦。
}

\item
%% number: 16
%% paperName: 2016年海南高考第 16 题 8 分
%% typeId: 6
%% type: 计算题
%% chapter: 光学
%% point: 全反射
%% method: 无
%% score: 8
%% degree: 450
%% duplicateId: 0
%% body:
{}[选修 3-4]（12 分）
\begin{enumerate}
	\item
	（4 分）下列说法正确的是\tkanswer[1.6cm]{ABD}。（填入正确答案标号。选对 1 个得 2 分，选对 2 个得 3 分，选对 3 个得 4 分；有选错的得 0 分）
	\fivechoices[answer=ABD]
	{在同一地点，单摆做简谐振动的周期的平方与其摆长成正比}
	{弹簧振子做简谐振动时，振动系统的势能与动能之和保持不变}
	{在同一地点，当摆长不变时，摆球质量越大，单摆做简谐振动的周期越小}
	{系统做稳定的受迫振动时，系统振动的频率等于周期性驱动力的频率}
	{已知弹簧振子初始时刻的位置及其振动周期，就可知振子在任意时刻运动速度的方向}
	\item
	（8 分）如图，半径为 $R$ 的半球形玻璃体置于水平桌面上，半球的上表面水平，球面与桌面相切于 A 点。一细束单色光经球心 O 从空气中摄入玻璃体内（入射面即纸面），入射角为 $45^{\circ}$，出射光线射在桌面上 B 点处。测得 AB 之间的距离为 $\frac{R}{2}$。现将入射光束在纸面内向左平移，求射入玻璃体的光线在球面上恰好发生全反射时，光束在上表面的入射点到 O 点的距离。不考虑光线在玻璃体内的多次反射。
	\onepicture[w=4.8cm, align=right]{\includesvg{figs/16}}
\end{enumerate}
%% answer:
\jdanswer{
\begin{enumerate}
	\item
	ABD。
	\item
	$\frac{\sqrt{2}}{2}R$。
\end{enumerate}
}
%% memo:
\memoanswer{
（1）由 $T=2\pi\sqrt{\frac{L}{g}}$，在同一地点，$g$ 相同，单摆周期的平方与其摆长成正比，故 A 正确；弹簧振子做简谐振动时，只有弹力做功，系统机械能守恒，故 B 正确；单摆运动周期与摆球质量无关，故 C 错误；受迫振动的频率等于驱动力的频率，与系统的固有频率无关，故 D 正确；由于不知道初始时刻的速度方向，无法写出振子的振动方程，故 E 错误。

（2）当光线经球心 O 入射时，光路图如图（a）所示。设玻璃的折射率为 $n$，由折射定律得 $n=\frac{\sin i}{\sin\gamma}$①，式中，入射角 $i=45^{\circ}$，$\gamma$ 为折射角，$\triangle OAB$ 为直角三角形，因此 $\sin\gamma=\frac{AB}{\sqrt{OA^{2}+AB^{2}}}$②。

\onepicture[w=5cm, align=right]{figs/16_an1.png}

发生全反射时，临界角 C 满足 $\sin C=\frac{1}{n}$③，在玻璃体球面上光线恰好发生全反射时，光路图如图（b）所示。设此时光线入射点为 E，折射光线射到玻璃体球面的 D 点。由题意有 $\angle EDO=C$④，在 $\triangle EDO$ 内，由正弦定理得 $\frac{OD}{\sin(90^{\circ}-\gamma)}=\frac{OE}{\sin C}$⑤。

\onepicture[w=5.2cm, align=right]{figs/16_an2.png}

联立以上各式并利用题给条件得 $OE=\frac{\sqrt{2}}{2}R$⑥。
}

\item
%% number: 17
%% paperName: 2016年海南高考第 17 题 8 分
%% typeId: 6
%% type: 计算题
%% chapter: 运动和力
%% point: 动量和能量
%% method: 无
%% score: 8
%% degree: 450
%% duplicateId: 0
%% body:
{}[选修 3-5]（12 分）
\begin{enumerate}
	\item
	（4 分）下列说法正确的是\tkanswer[1.6cm]{ACD}。（填入正确答案标号。选对 1 个得 2 分，选对 2 个得 3 分，选对 3 个得 4 分；有选错的得 0 分）
	\fivechoices[answer=ACD]
	{爱因斯坦在光的粒子性的基础上，建立了光电效应方程}
	{康普顿效应表明光子只具有能量，不具有动量}
	{玻尔的原子理论成功地解释了氢原子光谱的实验规律}
	{卢瑟福根据 $\alpha$ 粒子散射实验提出了原子的核式结构模型}
	{德布罗意指出微观粒子的动量越大，其对应的波长就越长}
	\item
	（8 分）如图，物块 A 通过一不可伸长的轻绳悬挂在天花板下，初始时静止；从发射器（图中未画出）射出的物块 B 沿水平方向与 A 相撞，碰撞后两者粘连在一起运动；碰撞前 B 的速度的大小 $v$ 及碰撞后 A 和 B 一起上升的高度 $h$ 均可由传感器（图中未画出）测得。某同学以 $h$ 为纵坐标，$v^{2}$ 为横坐标，利用实验数据作直线拟合，求得该直线的斜率为 $k=1.92\times10^{-3}\Usq/\Um$。已知物块 A 和 B 的质量分别为 $m_{A}=0.400\Ukg$ 和 $m_{B}=0.100\Ukg$，重力加速度大小 $g=9.80\Umsq$。
	\begin{enumerate}[label=(\roman*)]
		\item
		若碰撞时间极短且忽略空气阻力，求 $h-v^{2}$ 直线斜率的理论值 $k_{0}$；
		\item
		求 $k$ 值的相对误差 $\delta$（$\delta=\frac{\left|k-k_{0}\right|}{k_{0}}\times100\%$，结果保留 1 位有效数字）。
	\end{enumerate}
	\onepicture[w=3.2cm, align=right]{\includesvg{figs/17}}
\end{enumerate}
%% answer:
\jdanswer{
\begin{enumerate}
	\item
	ACD。
	\item
	（i）$2.04\times10^{-3}\Usq/\Um$；（ii）$6\%$。
\end{enumerate}
}
%% memo:
\memoanswer{
（1）爱因斯坦提出光子说，并建立了光电效应方程，故 A 正确；康普顿效应表明光子具有动量，故 B 错误；玻尔理论成功地解释了氢原子光谱，故 C 正确；在 $\alpha$ 粒子散射实验的基础上，卢瑟福提出了原子的核式结构模型，故 D 正确；由 $p=\frac{h}{\lambda}$，动量越大，其对应的波长越短，故 E 错误。

（2）（i）设物块 A 和 B 碰撞后共同运动的速度为 $v'$，由动量守恒定律得 $m_{B}v=(m_{A}+m_{B})v'$①，在碰后 A 和 B 共同上升的过程中，由机械能守恒定律得 $\frac{1}{2}(m_{A}+m_{B})v'^{2}=(m_{A}+m_{B})gh$②，联立①②式得 $h=\frac{m_{B}^{2}}{2g(m_{A}+m_{B})^{2}}v^{2}$③，由题意得 $k_{0}=\frac{m_{B}^{2}}{2g(m_{A}+m_{B})^{2}}$④，代入题给数据得 $k_{0}=2.04\times10^{-3}\Usq/\Um$⑤。

（ii）按照定义 $\delta=\frac{|k-k_{0}|}{k_{0}}\times100\%$⑥，由⑤⑥式和题给条件得 $\delta=6\%$⑦。
}

\end{enumerate}

\end{document}
